\documentclass[10pt,a4paper]{amsart}
\usepackage{amsmath,amssymb,amsthm,graphicx}
\usepackage{verbatim}
\newenvironment{reference}{\comment}{\endcomment}
%\newenvironment{reference}{}{}
\newenvironment{slogan}{\comment}{\endcomment}
\newenvironment{history}{\comment}{\endcomment}

% For commutative diagrams we use Xy-pic
\usepackage[all]{xy}

% We use 2cell for 2-commutative diagrams.
\xyoption{2cell}
% Theorem environments.
%
\theoremstyle{plain}
\newtheorem{theorem}[subsection]{Theorem}
\newtheorem{proposition}[subsection]{Proposition}
\newtheorem{lemma}[subsection]{Lemma}

\theoremstyle{definition}
\newtheorem{definition}[subsection]{Definition}
\newtheorem{example}[subsection]{Example}
\newtheorem{exercise}[subsection]{Exercise}
\newtheorem{situation}[subsection]{Situation}

\theoremstyle{remark}
\newtheorem{remark}[subsection]{Remark}
\newtheorem{remarks}[subsection]{Remarks}

\numberwithin{equation}{subsection}

% Macros
%
\def\lim{\mathop{\mathrm{lim}}\nolimits}
\def\colim{\mathop{\mathrm{colim}}\nolimits}
\def\Spec{\mathop{\mathrm{Spec}}}
\def\Hom{\mathop{\mathrm{Hom}}\nolimits}
\def\Ext{\mathop{\mathrm{Ext}}\nolimits}
\def\SheafHom{\mathop{\mathcal{H}\!\mathit{om}}\nolimits}
\def\SheafExt{\mathop{\mathcal{E}\!\mathit{xt}}\nolimits}
\def\Sch{\mathit{Sch}}
\def\Mor{\mathop{\mathrm{Mor}}\nolimits}
\def\Ob{\mathop{\mathrm{Ob}}\nolimits}
\def\Sh{\mathop{\mathit{Sh}}\nolimits}
\def\NL{\mathop{N\!L}\nolimits}
\def\CH{\mathop{\mathrm{CH}}\nolimits}
\def\proetale{{pro\text{-}\acute{e}tale}}
\def\etale{{\acute{e}tale}}
\def\QCoh{\mathit{QCoh}}
\def\Ker{\mathop{\mathrm{Ker}}}
\def\Im{\mathop{\mathrm{Im}}}
\def\Coker{\mathop{\mathrm{Coker}}}
\def\Coim{\mathop{\mathrm{Coim}}}

% Boxtimes
%
\DeclareMathSymbol{\boxtimes}{\mathbin}{AMSa}{"02}

%
% Macros for moduli stacks/spaces
%
\def\QCohstack{\mathcal{QC}\!\mathit{oh}}
\def\Cohstack{\mathcal{C}\!\mathit{oh}}
\def\Spacesstack{\mathcal{S}\!\mathit{paces}}
\def\Quotfunctor{\mathrm{Quot}}
\def\Hilbfunctor{\mathrm{Hilb}}
\def\Curvesstack{\mathcal{C}\!\mathit{urves}}
\def\Polarizedstack{\mathcal{P}\!\mathit{olarized}}
\def\Complexesstack{\mathcal{C}\!\mathit{omplexes}}
% \Pic is the operator that assigns to X its picard group, usage \Pic(X)
% \Picardstack_{X/B} denotes the Picard stack of X over B
% \Picardfunctor_{X/B} denotes the Picard functor of X over B
\def\Pic{\mathop{\mathrm{Pic}}\nolimits}
\def\Picardstack{\mathcal{P}\!\mathit{ic}}
\def\Picardfunctor{\mathrm{Pic}}
\def\Deformationcategory{\mathcal{D}\!\mathit{ef}}
\usepackage{cleveref}
\crefname{originallemma}{original lemma}{original lemmas}
\Crefname{originallemma}{Original lemma}{Original lemmas}
\makeatletter
\newlabel{algebra-lemma-flatness-descends-more-general}{{0584}{source}{}{sourcepointer-1}{}}
\newlabel{algebra-lemma-flatness-descends-more-general@cref}{{[originallemma][0][]0584}{[1][0][]source}{}{}{}}
\makeatother
\crefname{originallemma}{original lemma}{original lemmas}
\Crefname{originallemma}{Original lemma}{Original lemmas}
\makeatletter
\newlabel{lemma-flat-at-point}{{05VU}{source}{}{sourcepointer-2}{}}
\newlabel{lemma-flat-at-point@cref}{{[originallemma][0][]05VU}{[1][0][]source}{}{}{}}
\makeatother
\crefname{originallemma}{original lemma}{original lemmas}
\Crefname{originallemma}{Original lemma}{Original lemmas}
\makeatletter
\newlabel{lemma-flat-at-point-etale-local-rings}{{04NG}{source}{}{sourcepointer-3}{}}
\newlabel{lemma-flat-at-point-etale-local-rings@cref}{{[originallemma][0][]04NG}{[1][0][]source}{}{}{}}
\makeatother
\crefname{originallemma}{original lemma}{original lemmas}
\Crefname{originallemma}{Original lemma}{Original lemmas}
\makeatletter
\newlabel{spaces-properties-lemma-stalk-pullback-quasi-coherent}{{05VQ}{source}{}{sourcepointer-4}{}}
\newlabel{spaces-properties-lemma-stalk-pullback-quasi-coherent@cref}{{[originallemma][0][]05VQ}{[1][0][]source}{}{}{}}
\makeatother
\title{Flatness of a quasi-coherent module under a flat morphism of algebraic spaces}
\author{The Stacks Project Authors}
\date{stacks-project:spaces-morphisms.tex:local-snapshot-20261012; GNU FDL 1.2 or later}
\begin{document}
\maketitle
\setcounter{subsection}{1}
\noindent\textit{Editorial note.} All delivered mathematical text is the author's own, copied from the pinned Stacks Project chapter file spaces-morphisms.tex; no statement or proof step was rewritten, reordered or removed.\par\medskip
\noindent\textit{Editorial note.} The statement displays the pointwise equivalence and then states the global consequence; the proof establishes the displayed equivalence and the consequence follows by surjectivity.\par\medskip
\noindent\textit{Editorial note.} The four references of the counted spans are recorded as non-load-bearing navigation pointers; the argument that consumes them stays inside the counted proof span.\par\medskip
\noindent\textit{Editorial note.} Preamble: the three native excerpt spans of sources/stacks-spaces-morphisms/preamble.tex (lines 8-13, 14-18 and 154-217). The chapter file itself carries no preamble text.\par\medskip
\noindent\textit{Allowed standard prerequisites (editorial).}\begin{itemize}
\item Flatness at a point for modules on algebraic spaces. Lemmas lemma-flat-at-point and lemma-flat-at-point-etale-local-rings, used at proof lines 6524-6527. Reduce the sheaf statement to a statement about etale local rings.
\item The stalk of a pullback of a quasi-coherent sheaf. Properties of Spaces, Lemma spaces-properties-lemma-stalk-pullback-quasi-coherent, used at proof lines 6527-6529. Identifies the stalk of f\textasciicircum{}*G at the geometric point with a tensor product.
\item Flatness descends in the presence of the flat local ring map. Algebra, Lemma algebra-lemma-flatness-descends-more-general, used at proof lines 6535-6536. Is the algebraic statement the reduction yields.
\item Flatness at a point is tested on geometric points. used at proof lines 6521-6523. Allows the argument to be carried out at one geometric point.
\end{itemize}\medskip
\ifdefined\hypertarget\hypertarget{sourcepointer-1}{}\fi
\noindent\textit{Source reference.} Original lemma 0584: focused-seeds/AGP/sources/stacks-snapshot-20261011/tags.txt, source lines 6767--6767. Navigation pointer into the Stacks Project; the standard fact it supplies is applied inside the counted proof span and the argument resting on it stays there.\par\medskip
\ifdefined\hypertarget\hypertarget{sourcepointer-2}{}\fi
\noindent\textit{Source reference.} Original lemma 05VU: focused-seeds/AGP/sources/stacks-snapshot-20261011/tags.txt, source lines 7563--7563. Navigation pointer into the Stacks Project; the standard fact it supplies is applied inside the counted proof span and the argument resting on it stays there.\par\medskip
\ifdefined\hypertarget\hypertarget{sourcepointer-3}{}\fi
\noindent\textit{Source reference.} Original lemma 04NG: focused-seeds/AGP/sources/stacks-snapshot-20261011/tags.txt, source lines 6079--6079. Navigation pointer into the Stacks Project; the standard fact it supplies is applied inside the counted proof span and the argument resting on it stays there.\par\medskip
\ifdefined\hypertarget\hypertarget{sourcepointer-4}{}\fi
\noindent\textit{Source reference.} Original lemma 05VQ: focused-seeds/AGP/sources/stacks-snapshot-20261011/tags.txt, source lines 7559--7559. Navigation pointer into the Stacks Project; the standard fact it supplies is applied inside the counted proof span and the argument resting on it stays there.\par\medskip
\begin{lemma}
\label{lemma-flat-permanence}
Let $S$ be a scheme. Let $f : X \to Y$, $g : Y \to Z$ be morphisms of
algebraic spaces over $S$. Let $\mathcal{G}$ be a quasi-coherent sheaf on $Y$.
Let $x \in |X|$ with image $y \in |Y|$.
If $f$ is flat at $x$, then
$$
\mathcal{G}\text{ flat over }Z\text{ at }y
\Leftrightarrow
f^*\mathcal{G}\text{ flat over }Z\text{ at }x.
$$
In particular: If $f$ is surjective and flat, then
$\mathcal{G}$ is flat over $Z$, if and only if
$f^*\mathcal{G}$ is flat over $Z$.
\end{lemma}

\begin{proof}
Pick a geometric point $\overline{x}$ of $X$ and denote
$\overline{y}$ the image in $Y$ and $\overline{z}$ the image in $Z$.
Via the characterization of flatness in
Lemmas \ref{lemma-flat-at-point} and
\ref{lemma-flat-at-point-etale-local-rings}
and the description of the stalk of $f^*\mathcal{G}$ at $\overline{x}$ of
Properties of Spaces,
Lemma \ref{spaces-properties-lemma-stalk-pullback-quasi-coherent}
the lemma reduces to a purely algebraic question on the local
ring maps
$\mathcal{O}_{Z, \overline{z}} \to \mathcal{O}_{Y, \overline{y}}
\to \mathcal{O}_{X, \overline{x}}$
and the module $\mathcal{G}_{\overline{y}}$.
This algebraic statement is
Algebra, Lemma \ref{algebra-lemma-flatness-descends-more-general}.
\end{proof}

\end{document}
